\section{Conclusion}

The essay's conclusions fall into three categories, and they carry different evidentiary weight.

\textbf{Empirical finding.} The economic gradient of AI development is shifting. The cost of fixed capability is collapsing, and the premium for frontier capability decays quickly. Deployment is bifurcating: high-volume, low-stakes workloads migrate to commoditized open-weight models, while high-value, high-reliability workloads remain with frontier labs. Frontier labs grow in absolute dollars but lose relative share. The capability gradient is not closing uniformly---the benchmark gap remains around four months and has not consistently narrowed---but deployment migration is observable in gateway-traffic data and is consistent with broader deployment patterns. These are measured traffic patterns and inferred deployment trends, not claims about market-wide margins or pricing power, which the data cannot establish on their own. The empirical claims rest on industry production data (Vercel, Epoch AI) rather than peer-reviewed sources; the limitations of this evidence base are stated in Section~3 and should be treated as provisional pending independent verification.

\textbf{Theoretical interpretation.} The shift is consistent with a transition from product competition toward infrastructure competition---from margin extracted through differentiation to rent extracted through control over the conditions of access. Perez's installation/deployment distinction organizes this interpretation, but the interpretation is not demonstrated by the framework. The attractor is used heuristically: it denotes a family of stabilizing responses, not a state toward which the system is shown to evolve under specified dynamics. The indicators are proposed, not measured consistently over time.

\textbf{Unresolved causal hypothesis.} The shift may increase the incentive for frontier firms to favor governance designs that preserve their position in the infrastructure layer. This is one of three non-exclusive explanations, alongside genuine safety concern and regulatory uncertainty. The evidence reviewed here cannot adjudicate among them. The paper does not claim that the deployment transition caused the governance proposals; it claims that the proposals are consistent with the strategic-incentive hypothesis, and that the design audit stands regardless of which explanation is correct.

\textbf{Returning to the three hypotheses.} The framework's validation against the EU, China, and UK regimes (Sections~5.2--5.4) bears on H1, H2, and H3, though it cannot adjudicate among them. The design audit establishes that the design features are consequential under all three. The EU and UK regimes demonstrate that contestable designs are achievable---shared infrastructure, taxpayer funding, proportionality mechanisms, and alternatives analysis are all in operation. The China regime demonstrates that state-controlled designs can be contestable against firm capture while remaining opaque on publication. The Anthropic template's design signature---bespoke, firm-funded for the commercial evaluator, silent on open-weight---matches the prediction of H1 (compliance costs that fall on entrants), is also consistent with H3 (a federal standard that firms can shape), and does not exclude H2 (the safety concern may be genuine). The framework's contribution is to make the design signature visible as a choice, not to adjudicate the motive. If the draft bill, when introduced, adds government funding, shared infrastructure, proportionality provisions, and an open-weight regime, the design signature converges toward the EU and UK models and the exclusionary prediction is disconfirmed on those dimensions---regardless of which hypothesis best describes the intent.

\textbf{What remains uncertain.} Three things. First, the causal ordering between the economic gradient and the institutional response is not established. The safety proposals may be a response to genuine incidents rather than to the economic gradient; both may be true. Second, the codification pathway is contingent. The draft bill may fail, may be weakened, or may be strengthened by competing proposals. If it passes with government funding, shared infrastructure, proportionality provisions, and an open-weight regime, the exclusionary prediction is disconfirmed on those dimensions. Third, the layer analysis is incomplete. Whether a contestable layer survives in AI depends on policy choices not yet made---interoperability mandates, data portability, public compute---that are outside the scope of this paper.

\textbf{Analytical contribution.} The eight-dimension audit is the paper's independent contribution. It evaluates governance design without requiring a judgment about motive, and it is falsifiable per dimension. Applied to the Anthropic/METR arrangement, it finds the arrangement contestable on publication and potentially exclusionary on access, proportionality, and open-weight treatment---with the caveat that ``potentially'' marks where the burden of proof lies, not where the verdict is settled. Applied to the EU AI Act's conformity-assessment regime, the UK AI Safety Institute, and China's algorithm registration regime, the framework produces intelligible and discriminating results, identifying the EU and UK regimes as contestable on most dimensions and the China regime as contestable against firm capture but uncalibrated against state capture. The framework's value is that it makes the distributional consequences of design choices visible before codification fixes them, and that it compares design signatures across regimes.

\textbf{Research agenda.} Four questions follow. First, what is the distribution of control over AI infrastructure---compute, data centers, power, model weights---and how is it changing? Mapping the AI stack by ownership and access is a precondition for any governance analysis. Second, what design features distinguish safety-enhancing evaluation regimes from competition-distorting ones? The eight-dimension framework is a first pass; this paper's validation against the EU, China, and UK regimes suggests it is robust enough to extend but should add a state-capture dimension to handle regimes in which the state is the principal regulator. Third, what is the effect of compliance requirements on entry and competition? The fixed-cost argument is plausible but not quantified; a systematic comparison of compliance costs across jurisdictions and regimes would test whether the exclusionary prediction holds. Fourth, what evidence would support or challenge the phase-transition interpretation? The Perez framework generates expectations about investment priorities, innovation patterns, institutional arrangements, and the distribution of returns. These are testable; the paper has specified indicators in Section~2, and the next step is to measure them.

\textbf{Final paragraph.} Technological revolutions do not fail because they slow down. They fail because the institutional transition from installation to deployment is mismanaged. In the cases examined here, incumbent firms played a significant role in designing the governance arrangements that would apply to them. Whether this pattern generalizes to the wider policy process is a question for further research. What the paper establishes is narrower: that the design features of incumbent-participating governance can produce distributional effects independently of motive, and that the framework for evaluating those features is available now, before codification fixes them. The strongest version of the argument is not that AI safety governance is regulatory capture. It is that legitimate safety governance must be evaluated not only by the risks it mitigates, but also by the distribution of access, authority, and compliance costs it creates. That argument remains valuable even if the strategic-incentive hypothesis is never confirmed.
